% =====================================================================
%  ANNEXE E — RESSOURCES
% =====================================================================
\chapter{Ressources : quoi utiliser, et surtout quoi ne pas utiliser}

Le plus grand danger de l'apprenant moderne n'est pas le manque de ressources : c'est leur
abondance. Passer trois heures à comparer dix applications, c'est trois heures où vous ne
parlez pas. Cette annexe vous donne une sélection courte et justifiée --- puis une règle :
\textbf{choisissez une ressource par besoin, et gardez-la douze mois.}

\begin{piege}[Le piège du collectionneur]
Symptômes : vous avez installé sept applications, ouvert trois chaînes YouTube, créé cinq
comptes, téléchargé deux méthodes\ldots{} et votre niveau n'a pas bougé.\\
Remède : une application de vocabulaire, une source d'écoute, un partenaire de parole,
un dictionnaire. \textbf{Quatre outils, pas quarante.} La variété ne fait pas progresser ;
la répétition, si.
\end{piege}

\section{Le kit minimal : quatre outils, zéro franc}
\begin{tableau}[Ce dont vous avez réellement besoin]
\begin{tabularx}{\linewidth}{@{}p{3.0cm}p{4.0cm}X@{}}
\entete{Besoin} & \entete{Outil} & \entete{Comment l'utiliser} \tabularnewline
\hline
\textbf{Vocabulaire} & Anki (gratuit sur ordinateur et Android) &
15 cartes par jour, créées \textbf{par vous}, avec une phrase entière au recto. Révision quotidienne de 10 minutes. \tabularnewline
\textbf{Dictionnaire} & WordReference ou Cambridge Dictionary (gratuits) &
Cherchez un mot, \textbf{lisez la phrase d'exemple}, notez-la. Ne cherchez jamais plus de cinq mots à la suite. \tabularnewline
\textbf{Écoute} & Un podcast ou une chaîne pour apprenants (gratuit) &
Un épisode de 5 à 10 minutes, trois fois, selon le protocole des 5 passes du chapitre 16. \tabularnewline
\textbf{Parole} & Un humain : ami, collègue, club, application d'échange &
Trois fois par semaine, 20 minutes minimum. Rien ne remplace cette ligne du tableau. \tabularnewline
\end{tabularx}
\end{tableau}

\section{Dictionnaires et outils de référence}
\begin{tableau}[Les six références utiles]
\footnotesize
\begin{longtable}{@{}p{4.6cm}p{9.6cm}@{}}
\entete{Outil} & \entete{À quoi il sert} \tabularnewline
\endfirsthead
\entete{Outil} & \entete{À quoi il sert} \tabularnewline
\endhead
\en{WordReference} & Dictionnaire anglais-français très riche, avec forums où des milliers de nuances sont discutées. Idéal au début. \tabularnewline
\en{Cambridge Dictionary} & Dictionnaire anglais-anglais pour apprenants, avec transcription phonétique britannique \textbf{et} américaine, et audio. À privilégier dès le niveau A2. \tabularnewline
\en{Oxford Learner's Dictionaries} & Même usage, très fiable, avec l'\en{Oxford 3000} : la liste des 3\,000 mots essentiels --- exactement le volume visé par ce livre. \tabularnewline
\en{Linguee / DeepL Write} & Traducteur de \textbf{contextes} : vous voyez comment un mot est réellement employé dans des documents bilingues. Excellent pour vérifier une tournure. \tabularnewline
\en{Forvo} & Prononciation enregistrée par des natifs du monde entier. Indispensable pour un mot dont vous doutez. \tabularnewline
\en{YouGlish} & Tapez un mot ou une expression : le site trouve des vidéos YouTube où des natifs la prononcent en contexte réel. L'outil le plus utile pour l'accent. \tabularnewline
\end{longtable}
\end{tableau}

\begin{astuce}[Le réflexe anglais-anglais]
Dès le troisième mois, cessez de chercher \en{book} = « livre ». Cherchez la
\textbf{définition anglaise} et la \textbf{phrase d'exemple}. Cela double le temps de
consultation mais triple la mémorisation : vous apprenez le mot dans son habitat naturel,
et vous lisez de l'anglais au passage.
\end{astuce}

\section{Écoute : où trouver de l'anglais adapté}
\begin{tableau}[Podcasts et chaînes, par niveau]
\footnotesize
\begin{longtable}{@{}p{2.4cm}p{4.6cm}p{7.2cm}@{}}
\entete{Niveau} & \entete{Type de ressource} & \entete{Ce qu'il faut chercher} \tabularnewline
\endfirsthead
\entete{Niveau} & \entete{Type} & \entete{Critères de choix} \tabularnewline
\endhead
\textbf{A1--A2} & Podcasts pour apprenants (débit lent, vocabulaire contrôlé) &
Épisodes de 3 à 8 minutes, texte disponible. Exemples de formats à chercher : « 6 Minute English », « Learning English Broadcast », « English for beginners ». \tabularnewline
\textbf{A2--B1} & Chaînes pédagogiques, histoires racontées, actualités simplifiées &
Épisodes de 10 minutes avec transcript. Les sites d'actualités à trois niveaux (le même article en facile, moyen, difficile) sont parfaits. \tabularnewline
\textbf{B1--B2} & Podcasts conversationnels pour apprenants avancés, séries, interviews &
Débit naturel mais articulé ; sujets variés. Commencez \textbf{sans} sous-titres, puis vérifiez avec. \tabularnewline
\textbf{B2+} & Podcasts natifs : actualité, sciences, sport, humour, débats &
Ce que les Anglophones écoutent eux-mêmes. Objectif : comprendre sans aucun support. \tabularnewline
\end{longtable}
\end{tableau}

\begin{regle}[Comment choisir sa source d'écoute]
Trois critères, dans cet ordre : \textbf{(1) un transcript existe} (sinon le protocole des
5 passes est impossible) ; \textbf{(2) le sujet vous intéresse vraiment} --- on supporte
mieux la difficulté quand le contenu plaît ; \textbf{(3) une voix que vous aimez entendre}.
Vous écouterez cette personne 150 heures cette année : choisissez bien.
\end{regle}

\section{Lecture : des textes à votre taille}
\begin{tableau}[Progression de lecture sur 12 mois]
\begin{tabularx}{\linewidth}{@{}p{2.6cm}p{3.6cm}X@{}}
\entete{Période} & \entete{Type de texte} & \entete{Volume conseillé} \tabularnewline
\hline
Mois 1--3 & Textes écrits pour les apprenants, livres gradués niveau 1-2, articles très courts & 10 minutes par jour, 300 mots par texte. \tabularnewline
Mois 4--6 & Livres gradués niveau 3-4, articles d'actualité simplifiés, blogs personnels simples & 15 minutes par jour, un chapitre par semaine. \tabularnewline
Mois 7--9 & Romans jeunesse, premières pages de Wikipédia en anglais simple, articles de presse sur des sujets connus & 20 minutes par jour, 30 pages par semaine. \tabularnewline
Mois 10--12 & Romans pour adultes, presse générale, essais, blogs spécialisés & 30 minutes par jour, un livre par mois. \tabularnewline
\end{tabularx}
\end{tableau}

\begin{astuce}[Où trouver des livres gratuits et légaux]
Le projet Gutenberg (gutenberg.org) propose des dizaines de milliers de livres tombés dans
le domaine public, en anglais, téléchargeables en une fois. Commencez par des textes courts :
nouvelles, contes, récits de voyage. La version simplifiée de Wikipédia en anglais
(simple.wikipedia.org) est une mine : des milliers d'articles écrits en anglais facile.
Enfin, les bibliothèques municipales et les centres culturels étrangers prêtent souvent des
livres en anglais --- et organisent des clubs de conversation gratuits.
\end{astuce}

\section{Parler : trouver des partenaires}
\begin{tableau}[Les six voies possibles, de la plus accessible à la plus efficace]
\begin{tabularx}{\linewidth}{@{}p{3.2cm}p{2.0cm}X@{}}
\entete{Voie} & \entete{Coût} & \entete{Comment faire} \tabularnewline
\hline
\textbf{Se parler à soi-même} & 0 &
Narration continue de vos gestes, monologue sous la douche, réponses aux 30 questions du chapitre 15. C'est la voie la plus sous-estimée : elle crée l'automatisme sans peur du jugement. \tabularnewline
\textbf{Applications d'échange linguistique} & 0 &
Des applications de type « tandem » mettent en relation des francophones et des anglophones : 15 minutes en anglais, 15 minutes en français. Cherchez un partenaire sérieux, fixez un rendez-vous hebdomadaire, gardez-le un an. \tabularnewline
\textbf{Clubs et associations} & 0 à faible &
Clubs d'anglais universitaires, associations culturelles, centres culturels britanniques ou américains, groupes de discussion en ligne. L'avantage : la régularité et le collectif. \tabularnewline
\textbf{Amis et collègues anglophones} & 0 &
Proposez explicitement : « Pendant nos pauses, parlons anglais 15 minutes. » Beaucoup acceptent, personne n'ose demander. \tabularnewline
\textbf{Réseaux sociaux et messagerie} & 0 &
Commentez en anglais sous des vidéos, participez à des forums sur vos centres d'intérêt, échangez des messages vocaux. L'écrit spontané est un excellent pont vers l'oral. \tabularnewline
\textbf{Tuteurs en ligne} & payant &
Plateformes de cours particuliers à la demande. Même deux heures par mois suffisent à débloquer l'oral et à faire corriger vos erreurs récurrentes. Cherchez un tuteur qui \textbf{vous fait parler} 80\,\% du temps, pas un professeur qui explique. \tabularnewline
\end{tabularx}
\end{tableau}

\begin{mission}[La phrase à envoyer aujourd'hui]
Copiez, adaptez et envoyez ce message à trois personnes anglophones de votre entourage ---
ou à trois contacts d'une application d'échange :\\[4pt]
\en{Hello! My name is \ldots{} and I'm from Senegal. I've been learning English for about a
year and I'd love to practise speaking with a real person. Would you be interested in a
language exchange? We could talk for 20 minutes twice a week --- half in English, half in
French. No pressure at all, just let me know if you'd like to try!}\\[4pt]
Sur trois messages, vous obtiendrez statistiquement une réponse enthousiaste. C'est le
meilleur retour sur investissement de toute cette annexe.
\end{mission}

\section{Écrire et se faire corriger}
\begin{tableau}[Outils d'écriture]
\begin{tabularx}{\linewidth}{@{}p{4.0cm}X@{}}
\entete{Outil} & \entete{Usage} \tabularnewline
\hline
\textbf{Correcteurs automatiques} & Utiles pour l'orthographe et les erreurs de base. Mais ne les laissez \textbf{pas} corriger à votre place : lisez la suggestion, comprenez-la, puis écrivez-la vous-même. Sinon vous ne progresserez jamais. \tabularnewline
\textbf{Sites de correction par des natifs} & Certains sites permettent de soumettre un texte et de recevoir des corrections commentées par des anglophones. Renvoyez l'ascenseur en corrigeant des textes français. \tabularnewline
\textbf{Le journal personnel} & L'outil le plus puissant et le moins cher : 60 invites au chapitre 17, une par jour, 100 mots. Personne ne vous juge, donc vous osez des phrases ambitieuses. \tabularnewline
\textbf{Les échanges écrits réels} & Messages, e-mails, commentaires, forums. Écrire à un humain motive infiniment plus qu'écrire un exercice. \tabularnewline
\end{tabularx}
\end{tableau}

\begin{astuce}[La règle des deux couleurs]
Quand un texte vous est corrigé, recopiez-le : en \textbf{noir} ce que vous aviez écrit,
en \textbf{rouge} chaque correction. Relisez cette version une semaine plus tard. Vous
constaterez que 80\,\% de vos erreurs sont toujours les mêmes --- ce sont elles qu'il faut
inscrire dans le plan « trois erreurs par semaine » du chapitre 18.
\end{astuce}

\section{Grammaire : trois livres suffisent}
\begin{tableau}[Les ouvrages de référence classiques]
\begin{tabularx}{\linewidth}{@{}p{5.0cm}X@{}}
\entete{Ouvrage} & \entete{Pourquoi} \tabularnewline
\hline
\textbf{Une grammaire pratique pour apprenants} (type \emph{English Grammar in Use}) &
Une page de règle, une page d'exercices, corrigés inclus. Le format le plus efficace jamais inventé. Prenez le niveau intermédiaire. \tabularnewline
\textbf{Un ouvrage d'usage pour les questions fines} (type \emph{Practical English Usage}) &
Organisé comme un dictionnaire des difficultés. À consulter quand vous butez sur une nuance --- pas à lire d'une traite. \tabularnewline
\textbf{Un dictionnaire anglais-anglais pour apprenants} &
Voir la section précédente. C'est le livre que vous ouvrirez le plus souvent de toute votre année. \tabularnewline
\end{tabularx}
\end{tableau}

\begin{piege}[N'achetez pas trois méthodes]
Une méthode de grammaire, un dictionnaire, et ce livre : c'est tout. Ajouter une deuxième
méthode crée le sentiment de progresser sans produire le moindre progrès, parce que vous
passez votre temps à recommencer au chapitre 1 d'un autre ouvrage.
\end{piege}

\section{Applications : ce qu'elles font bien, ce qu'elles ne font pas}
\begin{tableau}[Verdict honnête]
\begin{tabularx}{\linewidth}{@{}p{4.4cm}p{4.6cm}X@{}}
\entete{Type d'application} & \entete{Ce qu'elle fait bien} & \entete{Ce qu'elle ne fait pas} \tabularnewline
\hline
\textbf{Applications ludiques de type « cours du jour »} &
Créer une habitude quotidienne, entretenir la motivation, donner les 800 premiers mots. &
Elles ne vous feront \textbf{jamais} parler. Leurs phrases sont trop courtes, trop propres, trop faciles. Utiles les trois premiers mois, puis insuffisantes. \tabularnewline
\textbf{Répétition espacée (cartes mémoire)} &
Mémoriser 3\,000 mots de façon durable --- c'est l'outil le plus efficace qui existe pour le vocabulaire. &
Elles n'enseignent ni la prononciation, ni la syntaxe, ni la spontanéité. \tabularnewline
\textbf{Reconnaissance vocale} &
Vérifier qu'un son est correct, donner un retour immédiat sans honte. &
Elle comprend trop bien : elle devine votre intention. Un humain, lui, ne devine pas. \tabularnewline
\textbf{Traduction automatique} &
Dépanner, vérifier un doute, donner un point de départ. &
Elle vous empêche de faire l'effort qui construit la compétence. À utiliser \textbf{après} avoir essayé, jamais avant. \tabularnewline
\end{tabularx}
\end{tableau}

\section{Sans internet : huit ressources que vous avez déjà}
Si votre connexion est limitée, rien n'est perdu. Voici ce qui fonctionne hors ligne.

\begin{enumerate}
  \item \textbf{Votre téléphone.} Changez la langue du système en anglais. Vous apprendrez
  200 mots d'interface en trois semaines, sans effort : \en{settings, share, save, delete,
  search, download, mute, brightness}.
  \item \textbf{La radio.} Les stations internationales anglophones émettent en ondes
  courtes ou en FM dans de nombreuses régions d'Afrique de l'Ouest. Vingt minutes par jour
  suffisent --- même en bruit de fond, votre oreille se calibre.
  \item \textbf{La télévision.} Un film que vous connaissez déjà, en version originale avec
  sous-titres anglais. Vous connaissez l'histoire : votre cerveau est libre de se concentrer
  sur la langue.
  \item \textbf{Les fichiers téléchargés.} Téléchargez podcasts et vidéos quand vous avez du
  réseau, puis écoutez-les hors ligne --- dans le bus, en marchant, en cuisinant. C'est du
  temps gratuit.
  \item \textbf{Les livres papier.} Un seul roman anglais lu trois fois vaut mieux que dix
  romans jamais ouverts. Les bibliothèques, librairies d'occasion et centres culturels en
  fournissent gratuitement.
  \item \textbf{Les humains autour de vous.} Un collègue, un voisin, un commerçant, un
  étudiant en anglais : créez un club de deux personnes qui se parlent en anglais une heure
  par semaine. C'est la ressource la plus puissante de cette liste.
  \item \textbf{Les produits de consommation.} Lisez les étiquettes, les notices, les
  emballages, les modes d'emploi en anglais. Vocabulaire concret et immédiatement réutilisable.
  \item \textbf{Ce livre.} Les dix textes d'écoute du chapitre 16, les trois textes de
  lecture du chapitre 17, les 600 phrases de l'annexe B, les 17 dialogues du chapitre 14 :
  vous avez déjà 20 heures de matériel hors ligne, corrigé et gradué.
\end{enumerate}

\begin{tableau}[Budget-type sur douze mois]
\begin{tabularx}{\linewidth}{@{}p{5.4cm}p{2.6cm}X@{}}
\entete{Poste} & \entete{Coût} & \entete{Commentaire} \tabularnewline
\hline
Anki + dictionnaires en ligne & 0 & Gratuit, sans publicité, sans limite. \tabularnewline
Podcasts et vidéos & 0 & Illimité. Le seul coût : vos données mobiles. \tabularnewline
Livres gradués d'occasion & faible & Deux ou trois suffisent pour l'année. \tabularnewline
Partenaire d'échange & 0 & Vous payez en français, eux en anglais. \tabularnewline
Tuteur en ligne (facultatif) & variable & 2 heures par mois, à partir du mois 4. \tabularnewline
Examen blanc en ligne (facultatif) & 0 à faible & Au mois 6 et au mois 12, pour une mesure objective. \tabularnewline
\end{tabularx}
\end{tableau}

\begin{memo}[Révision express --- Annexe E]
Quatre outils, pas quarante : \textbf{Anki} pour les mots, \textbf{un dictionnaire
anglais-anglais} pour le sens, \textbf{un podcast avec transcript} pour l'oreille,
\textbf{un humain} pour la bouche. Tout le reste est du bruit. Et souvenez-vous : la
meilleure ressource du monde ne fait rien à votre place. Vingt minutes par jour, pendant
365 jours, valent mieux que n'importe quelle collection d'applications.
\end{memo}
